% !TeX root = part5.tex

\documentclass[../final.tex]{subfiles}

\begin{document}

% При отдельной сборке продолжаем нумерацию полного конспекта.
\ifSubfilesClassLoaded{\setcounter{chapter}{2}\setcounter{section}{10}}{}


% ============================================================
% Лекция 29.10
% ============================================================

% ============================================================
\section{Комптоновское рассеяние на скалярной частице}
% ============================================================

% ------------------------------------------------------------
\subsection{Вершина с реальным фотоном}
% ------------------------------------------------------------

Потенциал взаимодействия после интегрирования по частям
удобно записывать в виде

\begin{equation*}
    \hat V
    =
    -e\left(A^\mu\partial_\mu-\partial_\mu A^\mu\right)
    -e^2A_\mu A^\mu,
\end{equation*}

где первое слагаемое отвечает вершине с одним фотоном,
второе --- вершине с двумя фотонами
(см. вершины скалярной электродинамики в разделе~\ref{sec:free_em_field}).
Для $\varphi=Ne^{-ipx}$ и $A^\mu=\varepsilon^\mu e^{-ikx}$

\begin{equation*}
    T_{fi}
    =
    -i\int\varphi_f^*\hat V\varphi_i\,d^4x
    =
    ieN_iN_f\,
    \underbrace{
        \varepsilon^\mu\left(p_i+k+p_f\right)_\mu
    }_{\varepsilon^\mu j_\mu}
    \int e^{i(p_f-p_i-k)x}\,d^4x.
\end{equation*}

\remark
Поскольку $\varepsilon k=0$, слагаемое $k$ в скобке можно
добавлять или убирать: $\varepsilon(p_i+k+p_f)=\varepsilon(p_i+p_f)$.
Запись с $k$ удобна, когда $p_f=p_i+k$:
тогда $\varepsilon(p_i+p_f)=\varepsilon(2p_i+k)$.


% ------------------------------------------------------------
\subsection{Диаграммы}
% ------------------------------------------------------------

Рассмотрим процесс

\begin{equation*}
    \gamma(k_1)+\pi(p_1)\rightarrow\gamma(k_2)+\pi(p_2).
\end{equation*}

В порядке $e^2$ ему соответствуют три диаграммы.

\begin{rbox}[left=18pt,right=18pt,top=14pt,bottom=10pt,unbreakable]{[]}
    \centering

    \begin{tikzpicture}[
        >=Latex,
        every node/.style={text=rtext},
        photon/.style={
            draw=rtext,
            line width=0.85pt,
            decorate,
            decoration={snake, amplitude=1.4pt, segment length=5pt}
        },
        scalar/.style={draw=rc3, line width=1.15pt}
    ]
        % a) s-канал
        \begin{scope}[shift={(-5.2,0)}]
            \coordinate (A) at (-0.7,0);
            \coordinate (B) at (0.7,0);
            \draw[scalar, ->] (-1.9,-1.2) -- (-1.3,-0.6);
            \draw[scalar] (-1.3,-0.6) -- (A);
            \draw[scalar] (A) -- (B);
            \draw[scalar, ->] (B) -- (1.3,-0.6);
            \draw[scalar] (1.3,-0.6) -- (1.9,-1.2);
            \draw[photon] (-1.6,1.2) -- (A);
            \draw[photon] (B) -- (1.6,1.2);
            \fill[rc4] (A) circle (1.5pt);
            \fill[rc4] (B) circle (1.5pt);
            \node[font=\small] at (-1.85,1.45) {$k_1$};
            \node[font=\small] at (1.85,1.45) {$k_2$};
            \node[font=\small] at (-2.1,-1.4) {$p_1$};
            \node[font=\small] at (2.1,-1.4) {$p_2$};
            \node[font=\small] at (0,-2.0) {а)};
        \end{scope}

        % b) u-канал
        \begin{scope}[shift={(0,0)}]
            \coordinate (A) at (-0.7,0);
            \coordinate (B) at (0.7,0);
            \draw[scalar, ->] (-1.9,-1.2) -- (-1.3,-0.6);
            \draw[scalar] (-1.3,-0.6) -- (A);
            \draw[scalar] (A) -- (B);
            \draw[scalar, ->] (B) -- (1.3,-0.6);
            \draw[scalar] (1.3,-0.6) -- (1.9,-1.2);
            \draw[photon] (-1.6,1.2) -- (B);
            \draw[photon] (A) -- (1.6,1.2);
            \fill[rc4] (A) circle (1.5pt);
            \fill[rc4] (B) circle (1.5pt);
            \node[font=\small] at (-1.85,1.45) {$k_1$};
            \node[font=\small] at (1.85,1.45) {$k_2$};
            \node[font=\small] at (-2.1,-1.4) {$p_1$};
            \node[font=\small] at (2.1,-1.4) {$p_2$};
            \node[font=\small] at (0,-2.0) {б)};
        \end{scope}

        % c) контактная
        \begin{scope}[shift={(5.0,0)}]
            \coordinate (C) at (0,0);
            \draw[scalar, ->] (-1.3,-1.2) -- (-0.65,-0.6);
            \draw[scalar] (-0.65,-0.6) -- (C);
            \draw[scalar, ->] (C) -- (0.65,-0.6);
            \draw[scalar] (0.65,-0.6) -- (1.3,-1.2);
            \draw[photon] (-1.2,1.2) -- (C);
            \draw[photon] (C) -- (1.2,1.2);
            \fill[rc4] (C) circle (1.5pt);
            \node[font=\small] at (-1.45,1.45) {$k_1$};
            \node[font=\small] at (1.45,1.45) {$k_2$};
            \node[font=\small] at (-1.5,-1.4) {$p_1$};
            \node[font=\small] at (1.5,-1.4) {$p_2$};
            \node[font=\small] at (0,-2.0) {в)};
        \end{scope}
    \end{tikzpicture}

    \captionsetup{
        font=footnotesize,
        labelfont=bf,
        skip=5pt
    }

    \captionof{figure}{
        Комптоновское рассеяние на бесспиновой частице:
        а) сначала поглощается $k_1$, затем испускается $k_2$;
        б) сначала испускается $k_2$, затем поглощается $k_1$;
        в) контактная двухфотонная вершина.
    }
    \label{fig:scalar_compton}
\end{rbox}


% ------------------------------------------------------------
\subsection{Пропагатор и амплитуды}
% ------------------------------------------------------------

Во внутренней линии диаграммы а) частица несёт импульс
$q=p_1+k_1$, причём $q^2\neq m^2$ --- она виртуальна. Соединение
двух токов через такую линию даёт множитель

\begin{equation*}
    \frac{1}{-q^2+m^2}
\end{equation*}

--- \rconcept{пропагатор} скалярной частицы (сравните с $-1/q^2$
для фотона, у которого $m=0$). Тогда

\begin{equation*}
    \mathcal M_a
    =
    \left(j_1\varepsilon\right)
    \left(j_2\varepsilon\right)
    \frac{1}{m^2-q^2}
    =
    e^2\,
    \frac{
        \left[\left(2p_2+k_2\right)\varepsilon_2\right]
        \left[\left(2p_1+k_1\right)\varepsilon_1\right]
    }{
        \left(p_1+k_1\right)^2-m^2
    },
\end{equation*}

\begin{equation*}
    \mathcal M_b
    =
    e^2\,
    \frac{
        \left[\left(2p_2-k_1\right)\varepsilon_1\right]
        \left[\left(2p_1-k_2\right)\varepsilon_2\right]
    }{
        \left(p_1-k_2\right)^2-m^2
    },
\end{equation*}

\begin{equation*}
    \mathcal M_c
    =
    2e^2\left(\varepsilon_1\varepsilon_2\right).
\end{equation*}

\remark
В знаменателе $\mathcal M_b$ стоит $(p_1-k_2)^2$: во внутренней
линии диаграммы б) частица уже испустила $k_2$.
Множитель $2$ в $\mathcal M_c$ возникает потому, что в $A^2$
каждый из двух фотонов можно сопоставить любому из двух
множителей $A$. Знаки слагаемых зависят от соглашений; ниже
используется только $\mathcal M_c$.

\subsubsection{Лабораторная система}

Пусть мишень покоится: $p_1=(m,\mathbf 0)$. Выберем калибровку
$\varepsilon^0=0$, тогда $\varepsilon_1p_1=\varepsilon_2p_1=0$ и
$\varepsilon_ik_i=0$. Значит,

\begin{equation*}
    \left(2p_1+k_1\right)\varepsilon_1=0,
    \qquad
    \left(2p_1-k_2\right)\varepsilon_2=0,
\end{equation*}

т.\,е. $\mathcal M_a=\mathcal M_b=0$ и
\rresult{в лабораторной системе работает только контактная
диаграмма}:

\begin{equation*}
    |\mathcal M|^2=4e^4\left(\varepsilon_1\varepsilon_2\right)^2.
\end{equation*}


% ------------------------------------------------------------
\subsection{Сечение}
% ------------------------------------------------------------

Поток налетающих фотонов $j=1$, поэтому

\begin{equation*}
    d\sigma
    =
    \frac{dW}{j}
    =
    \frac{|\mathcal M|^2}{2E_1\,2k_1}\,
    \frac{d^3p_2}{(2\pi)^3\,2E_2}\,
    \frac{d^3k_2}{(2\pi)^3\,2k_2}\,
    (2\pi)^4\delta^4(p_f-p_i).
\end{equation*}

Интегрируем по $d^3p_2$ с помощью $\delta^3$, затем по $dk_2$
с помощью $\delta(k_1+m-k_2-E_2)$. Поскольку
$E_2^2=m^2+|\mathbf k_1-\mathbf k_2|^2$,

\begin{equation*}
    \frac{d(k_2+E_2)}{dk_2}
    =
    \frac{k_1+m-k_1\cos\theta}{E_2},
\end{equation*}

и в лабораторной системе ($E_1=m$)

\begin{equation*}
    d\sigma
    =
    \frac{|\mathcal M|^2}{16\pi^2}\,
    \frac{1}{4k_1m}\,
    \frac{k_2\,d\Omega}{k_1+m-k_1\cos\theta}.
\end{equation*}

Здесь $\theta$ --- угол рассеяния фотона, $m$ --- масса частицы,
\rnote{на которой происходит рассеяние}. Из законов сохранения
следует \rresult{формула Комптона}

\begin{equation*}
    k_2
    =
    \frac{mk_1}{m+k_1\left(1-\cos\theta\right)},
\end{equation*}

так что $k_1+m-k_1\cos\theta=mk_1/k_2$, и

\begin{equation*}
    \frac{d\sigma}{d\Omega}
    =
    \frac{e^4}{16\pi^2m^2}\,
    \frac{\left(\boldsymbol\varepsilon_1\boldsymbol\varepsilon_2\right)^2}
         {\left[1+\dfrac{k_1}{m}\left(1-\cos\theta\right)\right]^2}.
\end{equation*}


% ------------------------------------------------------------
\subsection{Неполяризованные фотоны}
% ------------------------------------------------------------

\rassumption{Для неполяризованных частиц} нужно усреднить по двум
поляризациям начального фотона и просуммировать по поляризациям
конечного:
$\left(\boldsymbol\varepsilon_1\boldsymbol\varepsilon_2\right)^2
\to\left\langle\left(\boldsymbol\varepsilon_1\boldsymbol\varepsilon_2\right)^2\right\rangle$.
Сумма по поперечным поляризациям конечного фотона равна
$1-(\boldsymbol\varepsilon_1\hat{\mathbf k}_2)^2$, а усреднение по
$\boldsymbol\varepsilon_1\perp\mathbf k_1$ даёт

\begin{equation*}
    \left\langle
        \left(\boldsymbol\varepsilon_1\boldsymbol\varepsilon_2\right)^2
    \right\rangle
    =
    \frac{1}{2}\sum_{\lambda,\lambda'}
    \left(\boldsymbol\varepsilon_\lambda\boldsymbol\varepsilon'_{\lambda'}\right)^2
    =
    \frac{1}{2}\left(2-\sin^2\theta\right)
    =
    \frac{1+\cos^2\theta}{2}.
\end{equation*}

Итого

\begin{equation*}
    \frac{d\sigma}{d\Omega}
    =
    \frac{e^4}{16\pi^2m^2}\,
    \frac{1+\cos^2\theta}{2\left[1+\dfrac{k_1}{m}\left(1-\cos\theta\right)\right]^2}
    =
    \frac{r^2}{2}\,
    \frac{1+\cos^2\theta}{\left[1+\dfrac{k_1}{m}\left(1-\cos\theta\right)\right]^2},
    \qquad
    r=\frac{\alpha}{m}.
\end{equation*}

\remark
На лекции множитель $\tfrac12$ от усреднения по начальной
поляризации не был выписан. Он нужен, чтобы в пределе
$k_1\ll m$ получить правильное \rresult{томсоновское сечение}:

\begin{equation*}
    \sigma
    =
    \frac{r^2}{2}\int\left(1+\cos^2\theta\right)d\Omega
    =
    \frac{r^2}{2}\cdot2\pi\cdot\frac{8}{3}
    =
    \frac{8\pi}{3}r^2
    =
    \sigma_{\mathrm{Th}},
\end{equation*}

которое для электрона равно $665\,\mathrm{mb}$
(см. раздел о фотоэффекте в первой главе).

\begin{rbox}[left=18pt,right=18pt,top=12pt,bottom=10pt,unbreakable]{[]}
    \centering

    \begin{tikzpicture}
        \begin{axis}[
            width=0.78\linewidth,
            height=6cm,
            xmin=0, xmax=180,
            ymin=0, ymax=1.45,
            ytick={0,0.5,1},
            xtick={0,30,60,90,120,150,180},
            xlabel={$\theta$, градусы},
            ylabel={$\dfrac{1}{r^2}\dfrac{d\sigma}{d\Omega}$},
            axis line style={rtext},
            tick style={rtext},
            every tick label/.style={text=rtext},
            label style={text=rtext},
            legend style={
                draw=none, fill=none, text=rtext,
                at={(0.97,0.97)}, anchor=north east, font=\small
            },
            legend cell align=left
        ]
            \addplot[rconceptcolor, thick, smooth, samples=91, domain=0:180]
            {0.5*(1+cos(x)^2)};
            \addlegendentry{$k_1\ll m$ (Томсон)}

            \addplot[rc3, thick, smooth, samples=91, domain=0:180]
            {0.5*(1+cos(x)^2)/(1+0.5*(1-cos(x)))^2};
            \addlegendentry{$k_1=0.5\,m$}

            \addplot[rresultcolor, thick, smooth, samples=91, domain=0:180]
            {0.5*(1+cos(x)^2)/(1+2*(1-cos(x)))^2};
            \addlegendentry{$k_1=2\,m$}
        \end{axis}
    \end{tikzpicture}

    \captionsetup{
        font=footnotesize,
        labelfont=bf,
        skip=5pt
    }

    \captionof{figure}{
        Угловое распределение рассеянных фотонов.
        С ростом энергии фотона рассеяние назад подавляется
        и распределение вытягивается вперёд.
    }
    \label{fig:scalar_compton_angular}
\end{rbox}


% ============================================================
\section{Частицы со спином: уравнение Дирака}
% ============================================================

% ------------------------------------------------------------
\subsection{Идея Дирака}
% ------------------------------------------------------------

Уравнение Клейна--Гордона второго порядка по времени, поэтому
его плотность $\rho$ не положительно определена. Дирак искал
уравнение \rassumption{первого порядка} и по времени,
и по координатам:

\begin{equation*}
    i\frac{\partial\psi}{\partial t}
    =
    \hat H\psi
    =
    \left(\boldsymbol\alpha\hat{\mathbf p}+\beta m\right)\psi.
\end{equation*}

Каждая компонента $\psi$ должна при этом удовлетворять
уравнению Клейна--Гордона:

\begin{equation*}
    \hat H^2\psi
    =
    \left(\hat{\mathbf p}^2+m^2\right)\psi.
\end{equation*}


% ------------------------------------------------------------
\subsection{Условия на матрицы}
% ------------------------------------------------------------

Возводя гамильтониан в квадрат,

\begin{equation*}
    \left(\alpha_ip_i+\beta m\right)
    \left(\alpha_jp_j+\beta m\right)\psi
    =
    \Big[
        \alpha_i^2p_i^2
        +\sum_{i<j}\left(\alpha_i\alpha_j+\alpha_j\alpha_i\right)p_ip_j
        +\left(\alpha_i\beta+\beta\alpha_i\right)p_im
        +\beta^2m^2
    \Big]\psi.
\end{equation*}

Чтобы получить $\mathbf p^2+m^2$, нужно

\begin{equation*}
    \alpha_i^2=1,
    \qquad
    \beta^2=1,
    \qquad
    \left\{\alpha_i,\alpha_j\right\}=0\;(i\neq j),
    \qquad
    \left\{\alpha_i,\beta\right\}=0.
\end{equation*}

Значит, $\alpha_i$ и $\beta$ --- \rnote{не числа, а матрицы}.
Их свойства:

\begin{itemize}
    \item \rconcept{эрмитовость} (гамильтониан эрмитов):
    $A=A^\dagger=A^{*T}$;

    \item собственные значения $\lambda=\pm1$: из
    $Av=\lambda v$ следует $AAv=v=\lambda^2v$;

    \item нулевой след: $\operatorname{Tr}\alpha_i=0$,
    $\operatorname{Tr}\beta=0$. Например,
    $\alpha_i=-\beta\alpha_i\beta$, откуда
    $\operatorname{Tr}\alpha_i=-\operatorname{Tr}(\alpha_i\beta^2)
    =-\operatorname{Tr}\alpha_i$.
\end{itemize}

Приведённая к диагональному виду матрица
$A\to\operatorname{diag}(\lambda_1,\lambda_2,\ldots)$ с $\lambda=\pm1$
и нулевым следом имеет \rresult{чётную размерность}.
В размерности $2$ есть только три антикоммутирующие матрицы
(матрицы Паули), а нужно четыре. Поэтому минимальная размерность
--- \rresult{$4\times4$}, и $\psi$ --- четырёхкомпонентный
\rconcept{биспинор}.


% ------------------------------------------------------------
\subsection{Представление Дирака}
% ------------------------------------------------------------

Стандартный выбор:

\begin{equation*}
    \boldsymbol\alpha
    =
    \begin{pmatrix}
        0 & \boldsymbol\sigma\\
        \boldsymbol\sigma & 0
    \end{pmatrix},
    \qquad
    \beta
    =
    \begin{pmatrix}
        I & 0\\
        0 & -I
    \end{pmatrix},
\end{equation*}

где $\boldsymbol\sigma$ --- \rconcept{матрицы Паули}:

\begin{equation*}
    \sigma_1
    =
    \begin{pmatrix}
        0 & 1\\
        1 & 0
    \end{pmatrix},
    \qquad
    \sigma_2
    =
    \begin{pmatrix}
        0 & -i\\
        i & 0
    \end{pmatrix},
    \qquad
    \sigma_3
    =
    \begin{pmatrix}
        1 & 0\\
        0 & -1
    \end{pmatrix}.
\end{equation*}


% ------------------------------------------------------------
\subsection{Ковариантная форма}
% ------------------------------------------------------------

Умножим уравнение на $\beta$ и введём \rconcept{гамма-матрицы}

\begin{equation*}
    \gamma^\mu=\left(\beta,\,\beta\boldsymbol\alpha\right).
\end{equation*}

Получаем \rconcept{уравнение Дирака}

\begin{equation*}
    \left(i\gamma^\mu\partial_\mu-m\right)\psi=0.
\end{equation*}

Свойства гамма-матриц:

\begin{equation*}
    \gamma^\mu\gamma^\nu+\gamma^\nu\gamma^\mu=2g^{\mu\nu}I,
\end{equation*}

\begin{equation*}
    \left(\gamma^0\right)^2=I,
    \quad
    \gamma^{0\dagger}=\gamma^0,
    \qquad
    \left(\gamma^i\right)^2=-I,
    \quad
    \gamma^{i\dagger}=-\gamma^i,
\end{equation*}

\begin{equation*}
    \gamma^0
    =
    \begin{pmatrix}
        I & 0\\
        0 & -I
    \end{pmatrix},
    \qquad
    \boldsymbol\gamma
    =
    \begin{pmatrix}
        0 & \boldsymbol\sigma\\
        -\boldsymbol\sigma & 0
    \end{pmatrix}.
\end{equation*}

Для краткости будем писать $\hat p=\gamma^\mu p_\mu$.


% ------------------------------------------------------------
\subsection{Сохраняющийся ток}
% ------------------------------------------------------------

Введём \rconcept{дираковски сопряжённый спинор}

\begin{equation*}
    \bar\psi=\psi^\dagger\gamma^0.
\end{equation*}

Он удовлетворяет уравнению
$i\,\partial_\mu\bar\psi\,\gamma^\mu+m\bar\psi=0$.
Умножим уравнение Дирака слева на $\bar\psi$, сопряжённое ---
справа на $\psi$, и сложим:

\begin{equation*}
    \bar\psi\gamma^\mu\partial_\mu\psi
    +
    \left(\partial_\mu\bar\psi\right)\gamma^\mu\psi
    =
    \partial_\mu\left(\bar\psi\gamma^\mu\psi\right)
    =
    0.
\end{equation*}

Сохраняющийся ток $j^\mu=\bar\psi\gamma^\mu\psi$, и его
временная компонента

\begin{equation*}
    \rho
    =
    \bar\psi\gamma^0\psi
    =
    \psi^\dagger\psi
    =
    \sum_{i=1}^{4}|\psi_i|^2
    \geqslant0.
\end{equation*}

\rresult{Плотность положительно определена} --- в отличие
от уравнения Клейна--Гордона.


% ------------------------------------------------------------
\subsection{Покоящаяся частица}
% ------------------------------------------------------------

Ищем решение в виде

\begin{equation*}
    \psi=N\,u(p)\,e^{-ipx}.
\end{equation*}

Для частицы в покое $\mathbf p=0$ остаётся

\begin{equation*}
    \left(\gamma^0E-mI\right)u=0
    \qquad\Longleftrightarrow\qquad
    \begin{pmatrix}
        E & & & \\
         & E & & \\
         & & -E & \\
         & & & -E
    \end{pmatrix}
    \begin{pmatrix}
        u_1\\ u_2\\ u_3\\ u_4
    \end{pmatrix}
    =
    m
    \begin{pmatrix}
        u_1\\ u_2\\ u_3\\ u_4
    \end{pmatrix}.
\end{equation*}

Четыре решения:

\begin{equation*}
    E_1=m\;\text{(спин вверх)},
    \qquad
    E_2=m\;\text{(спин вниз)},
    \qquad
    E_3=-m,
    \qquad
    E_4=-m.
\end{equation*}

Решения с $E<0$ \rnote{нефизичны} как состояния частицы ---
их отбрасывают, а соответствующие им состояния описывают
\rconcept{античастицы}.


% ------------------------------------------------------------
\subsection{Движущаяся частица}
% ------------------------------------------------------------

Для движущейся частицы уравнение
$\left(i\gamma^\mu\partial_\mu-m\right)\psi=0$ даёт

\begin{equation*}
    \left(\hat p-m\right)u(p)=0.
\end{equation*}

Заметим, что

\begin{equation*}
    \left(\gamma^\mu p_\mu-m\right)
    \left(\gamma^\nu p_\nu+m\right)u(0)
    =
    \left(p^2-m^2\right)u(0)
    =
    0.
\end{equation*}

Следовательно, \rresult{$u(p)=(\hat p+m)\,u(0)$ --- решение
уравнения Дирака}:

\begin{equation*}
    u(p)
    =
    \left[
        \begin{pmatrix}
            I & 0\\
            0 & -I
        \end{pmatrix}E
        -
        \begin{pmatrix}
            0 & \boldsymbol\sigma\\
            -\boldsymbol\sigma & 0
        \end{pmatrix}\mathbf p
        +
        mI
    \right]u(0).
\end{equation*}

Взяв $u(0)=\begin{pmatrix}\chi\\0\end{pmatrix}$, получаем решения
с положительной энергией:

\begin{equation*}
    u_{1,2}(p)
    =
    N
    \begin{pmatrix}
        \chi\\[4pt]
        \dfrac{\boldsymbol\sigma\mathbf p}{E+m}\,\chi
    \end{pmatrix},
    \qquad
    \chi_1
    =
    \begin{pmatrix}
        1\\0
    \end{pmatrix},
    \quad
    \chi_2
    =
    \begin{pmatrix}
        0\\1
    \end{pmatrix}.
\end{equation*}

Аналогично, из $u(0)=\begin{pmatrix}0\\\chi\end{pmatrix}$ получаются
решения с отрицательной энергией ($E<0$):

\begin{equation*}
    u_{3,4}(p)
    =
    N
    \begin{pmatrix}
        \dfrac{-\boldsymbol\sigma\mathbf p}{-E+m}\,\chi\\[8pt]
        \chi
    \end{pmatrix}.
\end{equation*}

При нормировке $u^\dagger u=2E$ (как в уравнении Клейна--Гордона)
$N=\sqrt{|E|+m}$.

Для античастиц удобно перейти к положительной энергии и
импульсу, т.\,е. записать решение в виде

\begin{equation*}
    \psi=v(p)\,e^{ipx}.
\end{equation*}

\end{document}
